(1) 用分离变量和四条零位移边界得 \(\phi_{mn}(x,y)=\sin(m\pi x/a)\sin(n\pi y/b)\)，\(m,n\ge1\)，相应 \(\omega_{mn}=c\pi\sqrt{m^2/a^2+n^2/b^2}\)。不同模式在面积积分意义下正交，允许将初值展开为独立本征模。
(2) 当 \(a=b=L\) 时 \(\omega_{12}=\omega_{21}=c\pi\sqrt5/L\)，这是方形膜的几何对称简并。两模任意实系数线性叠加仍是同一频率的正常模，其节点可随叠加系数改变，但时变因子相同。
(3) 令 \(\omega=c\pi\sqrt5/L\)，解为 \(u(x,y,t)=[A\phi_{12}(x,y)+B\phi_{21}(x,y)]\cos\omega t\)。利用 \(\int_0^L\!\int_0^L\phi_{mn}^2dxdy=L^2/4\) 且交叉模积分为零，能量 \(E=\rho_sL^2\omega^2(A^2+B^2)/8\) 为常量；动能与张力势能在周期中交换。
(4) 对矩形膜，\(\omega_{12}^2=c^2\pi^2(1/a^2+4/b^2)\)，\(\omega_{21}^2=c^2\pi^2(4/a^2+1/b^2)\)，故 \(\omega_{12}^2-\omega_{21}^2=3c^2\pi^2(1/b^2-1/a^2)\)。不再简并时各自时间因子为 \(\cos\omega_{mn}t\)，两模叠加的节点一般随时间移动；差频小才形成缓慢拍振。
